\documentclass{cumcmthesis}

\title{2026 年全国大学生数学建模竞赛论文题目（三号黑体居中）}

\begin{document}

\begin{abstract}
这是 2026 年全国大学生数学建模竞赛电子版匿名提交标准 \LaTeX{} 论文模板摘要部分。摘要是论文的核心组成部分，应当言简意赅地概括全文的研究背景、基本假设、建模思想与机理分析、核心算法与求解策略、主要计算结果与量化结论，以及灵敏度分析与稳健性评价。

针对\textbf{问题一}，建立了……模型，利用……算法进行精确求解，计算得出……；

针对\textbf{问题二}，针对……特征，构建了……规划模型，采用……算法求解，求得最优决策参数为……，相比基准方案提升了……；

针对\textbf{问题三}，考虑……约束条件，建立了多目标/时空协同优化模型，通过……分析，得出了……；

针对\textbf{问题四/五}，构建了全局调度与对抗仿真模型，进行了全域优化，实现了……；

最后，对模型进行了离散收敛性检验、局部与全局灵敏度分析，并开展了 5000 次蒙特卡洛随机不确定性模拟，验证了模型在工程扰动下的高鲁棒性与实用价值。

\keywords{数学建模；时空协同优化；动态机理；差分进化算法；蒙特卡洛模拟}
\end{abstract}

\section{问题重述与背景分析}
\subsection{问题背景}
现代工程与国防应用中，针对动态目标防御与态势压制展开深入研究具有重要价值\cite{smith2018smoke, li2021uav}。

\subsection{问题重述}
本题需要依次解决以下几个核心问题：
\begin{enumerate}
    \item \textbf{问题 1}：建立基础动力学与观测几何模型，反演基准参数下的系统性能；
    \item \textbf{问题 2}：构建单执行机构最优控制与时序非线性规划模型并求解；
    \item \textbf{问题 3}：考虑动作间隔硬约束下的多阶段时序协同优化；
    \item \textbf{问题 4 与问题 5}：多智能体分布式协同汇聚与多目标对抗全局调度。
\end{enumerate}

\section{模型假设与符号说明}
\subsection{基本假设}
\begin{enumerate}
    \item 假设系统在任务周期内运行平稳，忽略外界微弱随机扰动；
    \item 假设各执行实体遵循经典动力学与刚体运动规律。
\end{enumerate}

\subsection{符号说明}
\begin{table}[H]
\centering
\caption{主要数学符号与物理量定义}
\begin{tabular}{ccl}
\toprule
\textbf{符号} & \textbf{单位} & \textbf{物理含义与说明} \\
\midrule
$t$ & $\text{s}$ & 系统时间变量 \\
$\bm{x}(t)$ & $\text{m}$ & 空间位置状态矢量 \\
$v$ & $\text{m/s}$ & 巡航运动速度 \\
$\theta$ & $\text{deg}$ & 水平航向角 \\
$T_{\text{obj}}$ & $\text{s}$ & 目标优化性能指标 \\
\bottomrule
\end{tabular}
\end{table}

\section{机理模型建立与解析几何推导}
在三维直角坐标系下，空间运动轨迹可表示为：
\begin{equation}
\bm{x}(t) = \bm{x}_0 + \bm{v} t + \frac{1}{2} \bm{a} t^2
\label{eq:kinematics}
\end{equation}

动态几何投影距离判据为：
\begin{equation}
d_{\text{min}}(t) = \left\| \bm{S}(t) - \bm{Q}(t) \right\| \le R_s
\end{equation}

\section{各问题求解与计算结果分析}
\subsection{问题 1 求解与分析}
如公式 \eqref{eq:kinematics} 所示进行离散求解，计算结果如表 \ref{tab:res1} 所示。

\begin{table}[H]
\centering
\caption{问题 1 计算数值结果汇总}
\label{tab:res1}
\begin{tabular}{cccc}
\toprule
\textbf{参数项} & \textbf{解析数值} & \textbf{单位} & \textbf{说明} \\
\midrule
起始时刻 & 8.0565 & $\text{s}$ & 判定截断点 \\
终止时刻 & 9.4480 & $\text{s}$ & 脱离视线 \\
有效时长 & 1.3915 & $\text{s}$ & 最终结果 \\
\bottomrule
\end{tabular}
\end{table}

\section{模型检验、灵敏度与鲁棒性分析}
对核心决策参数进行 $\pm 10\%$ 扰动分析，结果表明系统具备良好的稳定性与工程容错能力。

\section{模型评价与推广}
\subsection{模型优缺点}
本模型机理严密、解析性强、求解响应速度快；未来可进一步拓展至复杂三维环境。

\newpage
\bibliographystyle{plain}
\bibliography{references}

\newpage
\appendix
\section{核心程序代码清单}
\begin{verbatim}
import numpy as np

def solve_optimization():
    print("Executing 2026 CUMCM High-Precision Optimization Solver...")
    return True
\end{verbatim}

\end{document}
